\scanpage{10}
% 1–5
$k$ anneau commutatif
\[
\begin{aligned}
 S_n&=k[[t_1,\ldots,t_n]]\qquad(n\geq0)\\
 X_n&=\mathrm{Specf}(S_n)\\
 X_*&=(X_n)_{n\geq0}\quad\text{schéma formel simplicial augmenté}\\
 S_*&=(S_n)_{n\geq0}\quad\text{anneau top. cosimplicial coaugmenté}
\end{aligned}
\]

% 6–7
On suppose donné, pour tout $n\inserted{\geq0}$, un sous-ensemble
\[
 A_n\subset S_n
\]

% 8–11
On pose $A_*=(A_n)_{n\geq0}$. Un morphisme $X_m\xrightarrow{p}X_n$, associé à un hom de $k$-algèbres topologiques $S_n\xrightarrow{p^*}S_m$, est dit un $A_*$-\uline{morphisme} si $p^*(t_i)=f_i\in S_m$ est $\in A_m$ pour tout $1\leq i\leq n$.

% 11, left
% slanted
\margin{\uline{NB} les $f_i$ sont tels que $\varepsilon(f_i)$ \uline{nilpotents}}

% 12–13
Donc cela signifie que les $\mathrm{pr}_i\circ p$ $(1\leq i\leq n)$ sont des $A_*$-morph.

On suppose :

% 14–17
% slanted
\margin{Identités}
(a) Les opérations simpliciales entre les $X_n$ sont des $A_*$-morphismes ($\Longleftrightarrow\ \forall n\geq0$ et $1\leq i\leq n$ on a $t_i\in A_n$ ; $\Longleftrightarrow$ les morphismes identiques des $X_n$ sont $\in A_n$)

% 18–22
% slanted
\margin{Composition}
(b) Le composé de deux $A_*$-morphismes est \unsure{itou} [$\Longleftrightarrow$ si $p:X_m\to X_n$ est un $A_*$-morphisme, alors $p^*(A_n)\subset A_m\ \Longleftrightarrow$ Si $f_1,\ldots,f_n\in A_m$, et $\varepsilon(f_1),\ldots,\varepsilon(f_n)$ nilpotents, alors pour tout $\varphi\in A_n$ \inserted{avec $\varepsilon(\varphi)$ nilpotent}, on a $\varphi(f_1,\ldots,f_n)\in A_m$]

% 22, left
% slanted
\margin{NB \deleted{\unsure{d'ailleurs}}\inserted{Ce} \deleted{\unsure{que ce soit|qui est}}\inserted{sera} vrai \inserted{grâce à b)}, (sans supposer $\varepsilon(\varphi)$ nilpotent \deleted{dès que}) dès que $\mathrm{Im}(A_0\to S_n)\subset A_n$ \deleted{si $a\in A_0$, on}}

% 23–29
% slanted
\margin{\unsure{Condition annelée|Condition annelé}}
(c) \inserted{$1\in\unsure{A_0|A_n}$, et} \deleted{Le morphisme ``différence'' $X_2=X_1\mathbin{\widehat{\times}}_k X_1\to X_1$ $((x,y)\mapsto x-y)$ et ``produit'' [$(x,y)\mapsto xy$] sont des $A_*$ morphismes [\,$\Leftrightarrow t_1-t_2$ et $t_1t_2$ sont $\in A_2$, \inserted{et $1\in A_0$} $\Longleftrightarrow$} \deleted{les}\inserted{les} $A_n$ sont \deleted{stables par \deleted{\unsure{sommes et}} \inserted{différences et} produits, \deleted{et par différences} \inserted{et}}\inserted{des} \deleted{sous-groupes additifs stables par produit] et $1\in S_0$ [} sous-anneaux unitaires des $S_n$]
% archive sheet number: 9
